% Template LaTeX para TCC de graduação da UNIPAMPA — V2
% Copyright (c) 2026 Charles Quevedo Carpes.
% Autoria das contribuições próprias, organização e adaptações do template.
% This work may be distributed and/or modified under the conditions of the
% LaTeX Project Public License, version 1.3c or (at your option) any later version.
% License: https://www.latex-project.org/lppl/lppl-1-3c.txt (see LICENSE).
% This work has the LPPL maintenance status `author-maintained'.
% The Current Maintainer of this work is Charles Quevedo Carpes.
% Consulte AUTORIA.md para escopo, créditos e relação dos arquivos.

\chapter{REFERENCIAL TEÓRICO}

O referencial teórico reúne e discute a literatura relacionada ao problema investigado. Todas as fontes citadas no texto devem aparecer nas referências, e a lista final deve conter somente os documentos efetivamente citados.

\section{Citações diretas e indiretas}

Citações diretas de até três linhas permanecem no corpo do parágrafo, entre aspas duplas. Por exemplo: \enquote{a pesquisa científica exige planejamento e sistematização} \cite[p.~25]{gil2022}.

A citação direta com mais de três linhas deve ser apresentada sem aspas, com recuo de 4 cm, fonte 10 e espaçamento simples:

\begin{citacaodireta}
A citação longa é destacada do texto e deve reproduzir fielmente o trecho consultado. A chamada autor-data inclui o sobrenome em letras maiúsculas e minúsculas, o ano e a página ou localização da fonte consultada \cite[p.~88]{hall2000}.
\end{citacaodireta}

Supressões são indicadas por \texttt{[...]}; acréscimos ou comentários, por colchetes; e destaques acrescentados pelo autor do TCC devem ser identificados com a expressão \enquote{grifo nosso} na chamada.

\section{Citação de citação}

A citação de citação deve ser utilizada somente quando o documento original não puder ser consultado. O pacote oferece os comandos \verb|\apud| e \verb|\apudonline|. Exemplo: \apud[p.~89]{lacerda1998}{mariano2005}.

\section{Notas de rodapé}

Notas explicativas devem ser utilizadas com moderação e não podem interromper a sequência lógica do texto.\footnote{Exemplo de nota explicativa em fonte 10, espaço simples e alinhamento suspenso.}
